\section{Reasoning}
\label{sec:reasoning}

\subsection{Architettura a Due Livelli}

Il sistema di reasoning adotta un'architettura a due livelli che combina un
reasoner OWL 2 RL standard con un reasoner CONSTRUCT personalizzato per i
costrutti non coperti dal profilo RL.

\subsubsection{Livello 1: owlrl (OWL 2 RL)}

\texttt{owlrl} \citep{owlrl} implementa la chiusura deduttiva per il profilo
OWL 2 RL \citep{w3c-owl2-rl}, materializzando le triple inferite direttamente
nel grafo RDF. I costrutti supportati includono:
\begin{itemize}[nosep]
  \item \texttt{rdfs:subClassOf} --- propagazione del tipo verso l'alto;
  \item \texttt{owl:someValuesFrom} --- inferenza da restrizioni esistenziali
    (es. se G \texttt{wonAward} A, allora G \texttt{rdf:type AwardWinningGame});
  \item \texttt{owl:hasValue} --- inferenza da valore specifico;
  \item \texttt{owl:intersectionOf} --- scomposizione di intersezioni;
  \item \texttt{owl:inverseOf} --- materializzazione di propriet\`a inverse;
  \item \texttt{owl:SymmetricProperty} e \texttt{owl:TransitiveProperty};
  \item \texttt{owl:propertyChainAxiom} --- catene di propriet\`a.
\end{itemize}

Sull'insieme di istanze di test (3 giochi + 2 developer + piattaforme), \texttt{owlrl} produce
un incremento da $\sim$997 a $\sim$2680 triple.

\subsubsection{Livello 2: construct\_reasoner.py (SPARQL CONSTRUCT)}

21 regole CONSTRUCT colmano le lacune di OWL 2 RL, operando su costrutti che
richiedono ragionamento al di fuori del profilo RL. La Tabella~\ref{tab:construct_rules}
elenca le regole.

\begin{longtable}{lp{6cm}l}
\caption{Regole CONSTRUCT del reasoner personalizzato}\label{tab:construct_rules}\\
\toprule
\textbf{\#} & \textbf{Descrizione} & \textbf{Cosa inferisce} \\
\midrule
\endfirsthead
\multicolumn{3}{c}{\tablename\ \thetable{} -- continua} \\
\toprule
\textbf{\#} & \textbf{Descrizione} & \textbf{Cosa inferisce} \\
\midrule
\endhead
\bottomrule
\endlastfoot

R1 & HighlyRatedGame da metacriticScore $\geq$ 85     & \texttt{?g a vg:HighlyRatedGame} \\
R2 & RecentGame da releaseDate $\geq$ 2020-01-01      & \texttt{?g a vg:RecentGame} \\
R3 & StandaloneGame: VideoGame senza sequelOf          & \texttt{?g a vg:StandaloneGame} (negation-as-failure) \\
R4 & SequelGame: VideoGame con sequelOf                & \texttt{?g a vg:SequelGame} \\
R5a--f& Classificazione per piattaforma (PC, Console, Mobile, Nintendo, PS, Xbox) & \texttt{?g a vg:PCGame}, ecc. \\
R6a--c& Classificazione per modalit\`a (Multiplayer, SinglePlayer, Coop) & \texttt{?g a vg:MultiplayerGame}, ecc. \\
R7 & Inverse hasSequel: se A sequelOf B $\rightarrow$ B hasSequel A & \texttt{?b vg:hasSequel ?a} \\
R8 & Inverse awardWonBy & \texttt{?a vg:awardWonBy ?g} \\
R9a--d& Inverse: genreOf, platformFor, engineUsedIn, modeInGame & completamento bidirezionale \\
R10 & sharedGenreWith (pairwise) & \texttt{?g1 vg:sharedGenreWith ?g2} \\
R11 & sharedEngineWith (pairwise) & \texttt{?g1 vg:sharedEngineWith ?g2} \\
\bottomrule
\end{longtable}

Solo R1--R3 coprono costrutti realmente esterni a OWL 2 RL. Le regole R4--R9
(classificazione per sequel, piattaforma e modalit\`a, completamento delle inverse)
ripetono inferenze che il livello 1 gi\`a produce: sono mantenute perch\'e il livello 2
possa essere eseguito anche da solo, senza dipendere dalla chiusura \texttt{owlrl}.
R10 e R11 aggiungono relazioni di comodo tra giochi con lo stesso genere o motore.

\subsection{Perch\'e Alcuni Costrutti Non Sono in OWL 2 RL}

Il profilo OWL 2 RL \`e deliberatamente limitato per garantire la trattabilit\`a
computazionale (polinomiale); la Sezione~\ref{sec:profilo_dl} ne discute i limiti in
dettaglio. Per le classi dell'ontologia la situazione \`e la seguente: due famiglie di
costrutti esulano dal profilo, mentre la terza vi rientra ma non produce classificazioni.

\begin{enumerate}[nosep]
  \item \textbf{Datatype restriction} (\texttt{HighlyRatedGame}, \texttt{RecentGame}).
    Le restrizioni su tipi di dato con \texttt{xsd:minInclusive} richiedono il
    confronto di valori literal, operazione non prevista dalle regole RL standard.
    Le regole R1 e R2 risolvono con un semplice filtro SPARQL (\texttt{FILTER(?s >= 85)});
  \item \textbf{complementOf puro} (\texttt{StandaloneGame}). OWL 2 RL supporta
    solo forme limitate di negazione. La regola R3 implementa la
    \textit{negation-as-failure} con \texttt{FILTER NOT EXISTS}, assumendo la
    completezza dei dati (closed-world assumption sulla propriet\`a \texttt{sequelOf});
  \item \textbf{maxCardinality} (\texttt{GameWithSingleDeveloper}, \texttt{ExclusiveRelease}).
    Una restrizione ${\le}1$ in posizione di superclasse \`e ammessa in OWL 2 RL
    (regola \texttt{cls-maxc2}: se un'istanza ha due valori, questi vengono identificati
    con \texttt{owl:sameAs}). Poich\'e le due classi hanno sole condizioni necessarie
    (\texttt{rdfs:subClassOf}), per\`o, nessun gioco viene classificato in esse: fungono
    da vincoli di integrit\`a, e nessuna regola CONSTRUCT \`e necessaria.
\end{enumerate}

\subsection{Esempio di Inferenza a Due Livelli}

L'insieme di istanze di test (Appendice~\ref{sec:app_abox}) permette di seguire il
ragionamento su un caso concreto. Il gioco \texttt{GameA} (Elden Ring) \`e asserito
come \texttt{VideoGame} con data di rilascio 2022-02-25, punteggio Metacritic 96,
sviluppatore \texttt{Dev1}, due modalit\`a (individui di tipo \texttt{SinglePlayerMode}
e \texttt{MultiplayerMode}), due piattaforme (\texttt{PS5}, di tipo
\texttt{PlayStationPlatform}, e \texttt{PC\_Windows}, di tipo \texttt{WindowsPlatform})
e un premio. Anche \texttt{GameC} \`e sviluppato da \texttt{Dev1}.
La Tabella~\ref{tab:inferenza} riporta le conclusioni, il motivo e il livello che le produce.

\begin{table}[h]
\centering
\small
\caption{Inferenze su \texttt{GameA}}
\label{tab:inferenza}
\begin{tabularx}{\textwidth}{l X c}
\toprule
\textbf{Conclusione} & \textbf{Motivo} & \textbf{Livello} \\
\midrule
\texttt{PlayStationGame}, \texttt{ConsoleGame} &
  \texttt{availableOn PS5}; \texttt{PS5} \`e \texttt{PlayStationPlatform} $\sqsubseteq$
  \texttt{ConsolePlatform}; definizioni con $\exists$\texttt{availableOn} & 1 \\
\texttt{PCGame} & \texttt{WindowsPlatform} $\sqsubseteq$ \texttt{PCPlatform} & 1 \\
\texttt{SinglePlayerGame}, \texttt{MultiplayerGame} & $\exists$\texttt{hasGameMode} verso le due modalit\`a & 1 \\
\texttt{AwardWinningGame} & $\exists$\texttt{wonAward.Award}; il premio \`e \texttt{IndustryAward} $\sqsubseteq$ \texttt{Award} & 1 \\
\texttt{HighlyRatedGame} & $96 \ge 85$ (datatype restriction) & 2 \\
\texttt{RecentGame} & $2022\text{-}02\text{-}25 \ge 2020\text{-}01\text{-}01$ (datatype restriction) & 2 \\
\texttt{StandaloneGame} & nessun \texttt{sequelOf} noto (negation-as-failure) & 2 \\
\midrule
\texttt{Dev1 developerOf GameA} & inversa di \texttt{developedBy} & 1 \\
\texttt{GameA sharesDeveloperWith GameC} & catena \texttt{developedBy} $\circ$ \texttt{developerOf} e simmetria & 1 \\
\bottomrule
\end{tabularx}
\end{table}

L'esempio mostra la divisione dei compiti: il livello 1 (\texttt{owlrl}) copre
gerarchie, restrizioni esistenziali, inverse e catene; il livello 2 (CONSTRUCT)
interviene solo dove servono confronti tra literal o negazione.

\subsection{Bug Critico: owl:sameAs Falsi e Soluzione}

\subsubsection{Il Problema}

Durante lo sviluppo \`e emerso un bug critico: \texttt{owlrl} inferiva
\texttt{owl:sameAs} falsi tra giochi distinti a causa degli assiomi di
identit\`a (\texttt{FunctionalProperty}, \texttt{InverseFunctionalProperty},
\texttt{hasKey}) dichiarati su \textbf{DatatypeProperty}.

\textbf{Meccanismo del bug}: nel profilo OWL 2 RL le uguaglianze tra individui
sono prodotte da tre regole:
\begin{itemize}[nosep]
  \item \texttt{prp-fp} (FunctionalProperty) --- se lo \emph{stesso soggetto} ha due
    valori per la propriet\`a, i due \emph{valori} vengono identificati;
  \item \texttt{prp-ifp} (InverseFunctionalProperty) --- se \emph{due soggetti}
    condividono un valore, i due \emph{soggetti} vengono identificati;
  \item \texttt{prp-key} (\texttt{hasKey}) --- due individui della classe con gli
    stessi valori di chiave vengono identificati.
\end{itemize}
\texttt{owlrl} tratta i literal come nodi del grafo alla stregua di individui.
Applicate a datatype property, \texttt{prp-ifp} e \texttt{prp-key} identificano
giochi distinti che condividono il valore di un identificativo, mentre \texttt{prp-fp}
produce uguaglianze tra literal diversi dello stesso gioco (es.\ due date di rilascio
regionali). Ogni \texttt{owl:sameAs} spurio propaga poi tipi e propriet\`a tra gli
individui unificati (regole \texttt{eq-rep-*}), falsando la classificazione.

\textbf{Esempio}: nel test delle istanze (Appendice~\ref{sec:app_abox}) due
\texttt{owl:sameAs} spuri collegavano \texttt{GameA} agli altri due giochi,
propagando loro l'asserzione \texttt{wonAward}: prima del fix tutti e tre i giochi
risultavano \texttt{AwardWinningGame}; dopo il fix solo \texttt{GameA}, l'unico con
\texttt{wonAward} asserito, viene classificato come tale.

\subsubsection{La Soluzione}

La soluzione adottata (implementata in \texttt{enrich\_owl.py}) si articola in
quattro passi:

\begingroup\sloppy
\begin{enumerate}[nosep]
  \item \textbf{Nessuna InverseFunctionalProperty su datatype property}. Oltre a
    innescare \texttt{prp-ifp}, questo costrutto non \`e ammesso in OWL 2 DL;
    l'identit\`a dei giochi \`e espressa esclusivamente tramite
    \texttt{owl:hasKey} (su \texttt{hasWikidataId} e \texttt{hasSteamAppId});
  \item \textbf{Sezione ``DL-only axioms'' separata}. I 16 assiomi rimanenti
    (13 \texttt{FunctionalProperty} su datatype property e 3 \texttt{owl:hasKey})
    sono dichiarati in una sezione dedicata in coda a \texttt{videogames.owl},
    con un commento esplicativo;
  \item \textbf{Strip automatico prima del reasoning}. La funzione
    \texttt{strip\_dl\_only\_axioms()} rimuove questi assiomi dalla copia di
    lavoro in memoria prima del reasoning, prevenendo le inferenze errate;
  \item \textbf{Restore opzionale}. Con il flag \texttt{-{}-keep-dl-axioms}, la
    funzione \texttt{restore\_dl\_only\_axioms()} ripristina gli assiomi nel file
    serializzato; per default l'output resta privo di assiomi DL-only e quindi
    sicuro per \texttt{owlrl}.
\end{enumerate}
\endgroup

Lo schema \texttt{videogames.owl} conserva sempre la sezione DL-only, preservandone
il valore dichiarativo per l'uso con reasoner DL completi (HermiT, Pellet).

\subsubsection{Impatto}

Dopo il fix: 0 \texttt{owl:sameAs} falsi su tutto l'insieme di istanze di test (2715 triple finali,
nessun gioco unificato erroneamente). Le 12 famiglie di regole Python dirette in
\texttt{enrich\_owl.py} (\texttt{\_materialise}) classificano correttamente tutti
i giochi senza produrre falsi positivi.

\subsection{Materialiser Mirato in enrich\_owl.py}

Oltre al reasoning owlrl + CONSTRUCT, \texttt{enrich\_owl.py} implementa un
materialiser Python diretto con 12 famiglie di regole che operano senza passare
per SPARQL, risultando pi\`u veloce di \texttt{owlrl} full su insiemi di istanze di grandi
dimensioni:

\begin{enumerate}[nosep]
  \item AwardWinningGame $\leftarrow$ game \texttt{wonAward} \_;
  \item FranchiseGame $\leftarrow$ game \texttt{belongsTo} \_;
  \item sharedFranchiseWith $\leftarrow$ belongsTo $\circ$ includes
  \item sharesDeveloperWith $\leftarrow$ developedBy $\circ$ developerOf
  \item sharesPublisherWith $\leftarrow$ publishedBy $\circ$ publisherOf
  \item MultiplayerGame / SinglePlayerGame / CoopGame $\leftarrow$ \texttt{hasGameMode}
  \item PCGame / ConsoleGame / MobileGame $\leftarrow$ \texttt{availableOn}
  \item NintendoGame / PlayStationGame / XboxGame $\leftarrow$ \texttt{availableOn}
  \item HighlyRatedGame $\leftarrow$ \texttt{metacriticScore} $\geq$ 85
  \item RecentGame $\leftarrow$ \texttt{releaseDate} $\geq$ 2020-01-01
  \item SequelGame $\leftarrow$ game \texttt{sequelOf} \_
  \item StandaloneGame $\leftarrow$ VideoGame senza \texttt{sequelOf}
\end{enumerate}

Le regole pairwise (3--5) sono protette dal cap \texttt{-{}-max-group} (default 50)
per evitare esplosioni combinatorie su franchise con migliaia di giochi.
